\scanpage{17}
\deleted{des autom. de $P$ qui invarient le diviseur \deleted{\unsure{sections|section}} relatif de degré $2$ qui est défini par les $2$ sections de $P$, \uline{et} induisent l'identité dessus -- qui est un schém. en groupes $G$ lisse sur $S$, de dim relative $1$, dont la comp. \deleted{connexe}\inserted{neutre} $G^0$ est commutative}

\deleted{3°)}
% archive sheet number: 16
